\section*{الفصل \ref{ch:b2:cell-differentiation} --- تمايز الخلايا: الخلية العضلية الهيكلية}

\begin{solution}{exo:b2:cell-differentiation:1}
\omterm{def:b2:cell-differentiation:differentiation}{التمايز}: التعبير الثابت عن مورّثات صنف خلوي واحد و
وظائفه. \omterm{def:b2:cell-differentiation:differentiation}{والتحدد}: الالتزام بمصير قبل أن
يُعبَّر عنه. والقدرة: مدى المصائر التي لا تزال مفتوحة. فاللاقحة
\omterm{prop:b2:asexual-reproduction:tissueculture}{كلية القدرة}؛ وخلية \omterm{def:b2:mammal-reproduction:implantation}{الكتلة الخلوية الداخلية} متعددة القدرة؛ والخلية التابعة
أحادية القدرة (عضل فقط) --- وهي \omterm{def:b2:cell-differentiation:differentiation}{خلية جذعية} مع ذلك؛ \omterm{def:b2:cell-differentiation:myogenesis}{والليفة العضلية}
متمايزة، بعد انقسامية، بلا قدرة.
\end{solution}

\begin{solution}{exo:b2:cell-differentiation:2}
خلية الأديم الجلدي العضلي (Pax3) $\to$ \omterm{def:b2:cell-differentiation:myogenesis}{أرومة عضلية} مستحَثة (\omterm{def:b2:cell-differentiation:myogenesis}{MyoD}، وMyf5)،
تنقسم ما دام FGF باقيًا $\to$ خروج من الدورة، والميوجينين $\to$ اصطفاف
واندماج في \omterm{def:b2:cell-differentiation:myogenesis}{أنبوب عضلي} $\to$ مورّثات العضل مشتعلة، واللييفات العضلية
مركّبة (MRF4) $\to$ ليفة ناضجة بنوى محيطية، مع
خلايا تابعة ساكنة (Pax7).
\end{solution}

\begin{solution}{exo:b2:cell-differentiation:3}
أن نواة خلية متمايزة تحفظ كل مورّثة لازمة
لصنع حيوان كامل، وأن سيتوبلازم البيضة يستطيع إعادة ضبط
برنامجها: \omterm{def:b2:cell-differentiation:differentiation}{فالتمايز} عكوس وليس فقدًا لدنا. ولم
يبيّن أن كل نواة متمايزة يمكن إعادة ضبطها (فمعظم
النقول أخفقت)، ولا كيف تعمل إعادة الضبط، ولا أن الخلية
المتمايزة نفسها تستطيع تغيير مصيرها في مكانها.
\end{solution}

\begin{solution}{exo:b2:cell-differentiation:4}
قرصا Z يحدان \omterm{def:b2:cell-differentiation:fibre}{القسيم العضلي}؛ وخيوط الأكتين الرفيعة مرساة
إليهما؛ وخيوط الميوزين الغليظة متمركزة على الخط M؛ والنطاق A
(الخيوط الغليظة) والنطاق I (الرفيعة وحدها). وعند التقصير تنزلق الخيوط
الرفيعة نحو الخط M ويتقارب قرصا Z أحدهما من
الآخر: فينضيق النطاق I والمنطقة H، ولا يتغير النطاق A --- وهو طول
الخيط الغليظ --- ولا يقصر أي من الخيطين.
\end{solution}

\begin{solution}{exo:b2:cell-differentiation:5}
$0.2/2.5\times 10^{-6} = \num{80000}$ \omterm{def:b2:cell-differentiation:fibre}{قسيم عضلي}؛ و$0.2/25\times 10^{-6}
= \num{8000}$ نواة؛ \omterm{def:b2:cell-differentiation:myogenesis}{وأرومة عضلية} لكل نواة، أي نحو \num{8000}
\omterm{def:b2:cell-differentiation:myogenesis}{أرومة عضلية} اندمجت (وتضيف الخلايا التابعة مزيدًا لاحقًا).
\end{solution}

\begin{solution}{exo:b2:cell-differentiation:6}
يستحث FGF بروتين Id الذي يرتبط بالبروتين الشريك الذي يحتاجه \omterm{def:b2:cell-differentiation:myogenesis}{MyoD} لتكوين
ثنائي رابط للدنا؛ فيكون \omterm{def:b2:cell-differentiation:myogenesis}{MyoD} حاضرًا لكنه خامل، وتظل الخلايا
تنقسم. وحين يُسحب FGF يتحلل Id، ويتثنى \omterm{def:b2:cell-differentiation:myogenesis}{MyoD}، ويُشعل
الميوجينين، وتخرج الخلايا من الدورة وتندمج. أما الأرومات العضلية
التي تصنع Id باستمرار فلا تندمج أبدًا: بل تنقسم بلا نهاية و
لا تكوّن عضلًا.
\end{solution}

\begin{solution}{exo:b2:cell-differentiation:7}
الشرط $0.4m = m^{2}/(1 + m^{2}) + 0.02$. وعند $m = 2.1$: الطرف الأيسر
$0.84$، والأيمن $4.41/5.41 + 0.02 = 0.835$: أي حالة مستقرة. وأما الحالة
الدنيا: فمن أجل $m$ الصغيرة، $0.4m \approx m^{2} + 0.02$، ومنه $m \approx
0.055$ (الأيسر $0.022$، والأيمن $0.023$). وإذا رُفع إلى 1، وهو فوق
العتبة قرب $0.41$، تسلقت الخلية إلى $2.1$ وبقيت: فقد
التزمت.
\end{solution}

\begin{solution}{exo:b2:cell-differentiation:8}
مع $k = 0.6$: عند $m = 1$، يفوق التحلل $0.6$ الاصطناع
$0.52$؛ وعند $m = 0.5$، $0.30 > 0.22$؛ وعند $m = 2$، $1.2 > 0.82$: أي إن
المستقيم يعلو المنحني في كل مكان وراء الحالة الدنيا، وهي
الحالة المستقرة الوحيدة. والخلية التي يتحلل \omterm{def:b2:cell-differentiation:myogenesis}{MyoD} فيها أسرع مما ينبغي
لا تستطيع مسك الحالة المشتعلة: فترتد إلى الحالة غير
المتمايزة مهما كان الاستحثاث --- أي لا عضل.
\end{solution}

\begin{solution}{exo:b2:cell-differentiation:9}
(أ) لا أرومات عضلية ولا عضل هيكلي: فالعاملان زائدان
أحدهما عن الآخر في \omterm{def:b2:cell-differentiation:differentiation}{التحدد}، فلا بد من زوالهما معًا. (ب) تتكون الأرومات العضلية وتصطف
لكنها لا تندمج أبدًا ولا تصنع بروتينات العضل: فالميوجينين لازم
\omterm{def:b2:cell-differentiation:differentiation}{للتمايز}. (ج) شبه سوي، لأن Myf5 يحل محل
\omterm{def:b2:cell-differentiation:myogenesis}{MyoD} في \omterm{def:b2:cell-differentiation:differentiation}{التحدد}.
\end{solution}

\begin{solution}{exo:b2:cell-differentiation:10}
في الأرومة الليفية تقع مورّثات العضل في صبغين يستطيع \omterm{def:b2:cell-differentiation:myogenesis}{MyoD}
فتحه؛ وفي الخلية الكبدية قد تكون محشوة في صبغين متغاير، أو قد
تكبتها منظّمات الكبد الرئيسة، فلا تكون الخلية
أهلًا. والاختبار: عالج خلايا كبدية بدواء يرخي
الصبغين (مثبط لنازعة أسيتيل الهستونات) قبل إضافة \omterm{def:b2:cell-differentiation:myogenesis}{MyoD}، أو
أضف \omterm{def:b2:cell-differentiation:myogenesis}{MyoD} مع عامل يزيح برنامج
الكبد، وقِس التعبير عن الميوزين.
\end{solution}

\begin{solution}{exo:b2:cell-differentiation:11}
الإطلاق المستمر المنخفض التواتر يرفع الكالسيوم داخل الخلية
مددًا طويلة؛ وتشعل المسالك المفعَّلة بالكالسيوم مورّثات
الميوزين البطيء وتكوّن المتقدرات والميوغلوبين ونمو
الشعيرات، وتطفئ الأشكال السريعة. أما برنامج الليفة
الرئيس (حالة \omterm{def:b2:cell-differentiation:myogenesis}{MyoD}، ومعمار \omterm{def:b2:cell-differentiation:fibre}{القسيم العضلي}) فلا يُمَسّ:
فالتدريب يغير أي مورّثات العضل يُعبَّر عنها، لا
هوية الخلية، ولا إشارة من عصب تستطيع إعادة فتح
برنامج الغضروف الذي أُغلق عند \omterm{def:b2:cell-differentiation:differentiation}{التحدد}.
\end{solution}

\begin{solution}{exo:b2:cell-differentiation:12}
الحالة المشتعلة لعامل ينشّط نفسه نقطة مستقرة من
حلقة تغذية راجعة، يبقيها اصطناع مستمر، لا أي تغير في
المورّثات؛ وتضيف علامات الصبغين طبقة ثانية تبقي المورّثات
المسكتة مغلقة. ولذلك تقتضي إعادة الضبط دفع الدارة
دون عتبتها وإعادة فتح الصبغين في آن، وهو ما يستطيعه
سيتوبلازم البيضة أو العوامل الأربعة في كسر صغير من
الخلايا فقط: فهي ممكنة، لأن لا شيء يضيع، وقليلة الكفاءة، لأن
لا بد من فتح قفلين مستقلين معًا.
\end{solution}

\begin{solution}{pb:b2:cell-differentiation:1}
\textbf{1.} $0.2/2.5\times 10^{-6} = \num{80000}$ \omterm{def:b2:cell-differentiation:fibre}{قسيم عضلي}؛
و$0.2/25\times 10^{-6} = \num{8000}$ نواة.
\textbf{2.} نحو \num{8000} في الليفة؛ و$4\times 10^{9}$ في
العضلة.
\textbf{3.} $\pi(30\times 10^{-6})^{2}\times 0.2 = \qty{5.7e-10}{m^3}$
في الليفة؛ و$2.8\times 10^{-4}$ \unit{m^3}، أي \qty{0.28}{L}، في
العضلة.
\textbf{4.} $5.7\times 10^{-10}/8000 = \qty{7e-14}{m^3} =
\qty{70000}{\micro m^3}$: أي ما يعادل 35 خلية عادية لكل نواة.
\textbf{5.} الداخل محشو بلييفات عضلية يجب أن تجري
متصلة من طرف إلى طرف؛ والنوى عند الحافة تترك الشبكة
متصلة وتجلس قرب الغشاء وإشاراته.
\textbf{6.} $4\times 10^{9}/2000 = 2\times 10^{6}$: أي 21 مضاعفة، و10.5
يوم.
\textbf{7.} $0.4m = m^{2}/(1 + m^{2}) + 0.02$. وعند $m = 0.055$: الأيسر
$0.022$، والأيمن $0.003 + 0.02 = 0.023$.
\textbf{8.} عند $m = 2.1$: الأيسر $0.84$، والأيمن $0.815 + 0.02 = 0.835$.
\textbf{9.} عند $m = 0.41$: الأيسر $0.164$، والأيمن $0.144 + 0.02 = 0.164$.
وتحته مباشرة يقل الاصطناع عن التحلل، فيهبط $m$
نحو $0.055$؛ وفوقه مباشرة يفوق الاصطناع التحلل ويرتفع $m$
نحو $2.1$: فأي إزاحة تكبر --- أي غير مستقرة.
\textbf{10.} $0.6 > 0.41$: فتذهب الخلية الأولى إلى الحالة المشتعلة و
تتمايز؛ و$0.3 < 0.41$: فترتخي الثانية إلى $0.055$ ولا
تتمايز.
\textbf{11.} مع $b = 0.4$ يبدأ منحني الاصطناع عند $0.4$ و
يبقى فوق المستقيم $0.4m$ حتى $m \approx 3.3$: فيزول التقاطع
الأدنى. فتتمايز الخلية تلقائيًا، بلا استحثاث.
\textbf{12.} الانقسام ينصّف $m$؛ فترث البنتان $2.1/2 =
1.05$، وهو فوق العتبة $0.41$، وتتسلقان عائدتين إلى $2.1$: فتُعاد
الحالة في كل بنت. ويبقى الالتزام ما دام
المستوى المنصَّف فوق العتبة.
\textbf{13.} $4$ في المليمتر $\times$ \qty{200}{mm} $= 800$ خلية تابعة
في الليفة؛ ويجب تعويض \qty{10}{\%} من \num{8000} نواة، أي \num{800}.
\textbf{14.} يحوي المقطع المصاب $80$ خلية تابعة؛ و$800/80 =
10$، أي 4 مضاعفات ($\times 16$، فتعطي \num{1280}: تندمج منها 800، وتبقى 480
لترميم المخزون) --- أي \qty{72}{h}، ثلاثة أيام.
\textbf{15.} ثلاثة أيام من الانقسام في مقابل أسبوعين إلى ثلاثة
مرصودة: والباقي إزالة الحطام بالبالعات، وتأخر
التنشيط، والاندماج، وتركيب اللييفات العضلية، وإعادة التعصيب و
نضج المقطع الجديد.
\textbf{16.} لا بد أن يعيد \omterm{def:b2:cell-differentiation:myogenesis}{MyoD} (والميوجينين) إشعال مورّثات العضل
في الخلايا المنشَّطة؛ والانقسام \omterm{def:b2:cell-differentiation:differentiation}{والتمايز} متنافيان
--- فلا بد أن يهبط Id وتتوقف الدورة قبل أن يعمل الميوجينين و
يقع الاندماج.
\textbf{17.} $0.95^{n} = 0.5$: ومنه $n = 13.5$، أي نحو أربع عشرة إصابة.
\textbf{18.} $800\times 0.95^{20} = 800\times 0.36 \approx 290$.
\textbf{19.} كل تقلص يمزق ليفات، وكل تمزق يستدعي
المخزون، وينكمش المخزون بضعة في المئة في كل مرة؛ فبعد سنوات
ينفد، ولا تُعاد بناء الليفات الممزقة، وتملأ الأرومات الليفية
الفجوات بالكولاجين.
\textbf{20.} $(80/60)^{2} = 1.78$؛ والنوى $8000\times 1.78 = \num{14200}$:
أي نحو \num{6200} نواة زائدة من الخلايا التابعة.
\textbf{21.} $0.28\times 1.78 = \qty{0.5}{L}$.
\textbf{22.} المتغير: مورّثات أشكال الميوزين، وكثافة المتقدرات
والشعيرات، والميوغلوبين. وغير المتغير: هوية العضل (حالة
\omterm{def:b2:cell-differentiation:myogenesis}{MyoD})، ومعمار \omterm{def:b2:cell-differentiation:fibre}{القسيم العضلي}، والنوى ومورثومها.
\textbf{23.} الليفات بعد انقسامية ولا تستطيع التكاثر؛ فالطريقان
الوحيدان ليفات أثخن ونوى أكثر في الليفة من الخلايا
التابعة.
\textbf{24.} التدريب: تثخن محدود، إذ لا تستطيع الليفة إضافة
نوى؛ والإصابة: لا إصلاح --- فتُعوَّض المقاطع المتلفة
بنسيج ليفي، كما في الحثل.
\textbf{25.} \num{8000} نواة في الليفة؛ والعتبة $m \approx 0.41$؛
ونحو ثلاثة أيام من انقسام الخلايا التابعة لإصلاح.
\end{solution}
